\section{Planejamento, autonomia e técnicas recentes aplicadas a operações de proximidade}
\label{sec_planejamento_autonomia_tecnicas_recentes}

Nesta seção é apresentada uma revisão da literatura relacionada ao planejamento, à autonomia e às técnicas recentes aplicadas a operações de proximidade orbital. São destacadas as referências que tratam do planejamento de trajetória, do guiamento sob restrições, da prevenção de colisão, da tomada de decisão autônoma e da incorporação de técnicas recentes, incluindo métodos baseados em aprendizado de máquina, em missões de \textit{rendezvous}, \textit{docking} e \textit{on-orbit servicing}.

Nos desenvolvimentos mais recentes da literatura, as operações de proximidade passaram a ser tratadas não apenas como problemas de rastreamento de trajetória ou de estabilização local, mas também como tarefas integradas de percepção, planejamento, tomada de decisão e controle sob restrições. Essa evolução está associada à crescente complexidade das missões de \textit{on-orbit servicing}, remoção ativa de detritos e \textit{rendezvous} com alvos não cooperativos, nas quais a espaçonave perseguidora deve operar em ambientes geometricamente restritos, com incertezas de estado, risco de colisão e necessidade de maior autonomia.

\cite{alizadeh2024survey} apresentam uma revisão sobre robótica espacial para \textit{on-orbit servicing}, destacando a interação entre estimação do estado do alvo, planejamento de movimento, controle em malha fechada e técnicas de aprendizado de máquina. Os autores mostram que o avanço recente da área depende cada vez mais da integração entre essas camadas funcionais, e não apenas do desenvolvimento isolado de subsistemas. Essa referência é importante por situar o problema do planejamento e da autonomia em um contexto mais amplo de operação robótica espacial.

\cite{dong2022safe} apresentam uma estrutura de controle preditivo voltada ao \textit{rendezvous} e \textit{docking} autônomos com alvo cambaleante em múltiplos estágios. Os autores consideram restrições práticas como saturação de atuadores, limites de velocidade e prevenção de colisão em condições adequadas ao acoplamento final. A formulação é construída de modo a garantir segurança e estabilidade, combinando etapas de aproximação global e refinamento terminal. Os resultados apresentados mostram que a abordagem conduz a espaçonave perseguidora a regiões compatíveis com o acoplamento, preservando requisitos de segurança operacional. Essa referência mostra uma tendência importante da literatura recente em direção a arquiteturas de planejamento e controle formuladas explicitamente sob restrições.

O problema de manobras de proximidade com controle preditivo é estudado por \cite{wang2023close} para \textit{rendezvous} e \textit{docking} com cliente rotativo ou não rotativo, modelando explicitamente a \textit{keep-out zone} do alvo por meio de uma elipsoide expandida. Os autores buscam melhorar simultaneamente o desempenho e a eficiência computacional do controlador, aspecto importante para aplicações em tempo real. Essa referência é relevante porque aproxima a formulação matemática das condições operacionais encontradas em missões reais, nas quais não basta convergir à vizinhança do alvo, sendo necessário respeitar zonas seguras, evitar colisões e manter viabilidade computacional ao longo de toda a aproximação.

Mais recentemente, \cite{madonna2025guidance} apresentam uma arquitetura de guiamento e controle para \textit{rendezvous} a curta distância e aproximação final de uma espaçonave perseguidora em direção a um alvo não cooperativo. O estudo combina planejamento de trajetória, lógica de guiamento e tratamento explícito da geometria da aproximação final, com foco em operações de curta distância. Essa referência mostra uma tendência mais recente da literatura, na qual a fase terminal da operação passa a ser tratada como uma etapa com requisitos geométricos, cinemáticos e operacionais próprios, e não apenas como um prolongamento da dinâmica relativa linearizada.

Paralelamente às abordagens baseadas em otimização e controle preditivo, a literatura recente também passou a explorar técnicas de aprendizado por reforço para ampliar a autonomia decisória em cenários com obstáculos, incertezas e necessidade de resposta rápida. \cite{sharma2024safety} apresentam uma metodologia de guiamento para \textit{rendezvous} baseada em \textit{safety reinforcement learning}, formulando o problema como um processo de decisão de Markov em situações que envolvem prevenção de colisão. Os autores mostram que a estratégia proposta é capaz de realizar desvio de obstáculos e atingir \textit{rendezvous} de alta precisão, ao mesmo tempo em que busca lidar com uma limitação recorrente dos métodos tradicionais de aprendizado por reforço, associada à dificuldade de impor restrições de segurança durante a exploração. Essa referência mostra que a autonomia em missões de proximidade passou a envolver não apenas automação do controle, mas também maior delegação do processo decisório.

No contexto da robótica espacial acoplada, \cite{alali2024reinforcement} investigam o planejamento de trajetória para um manipulador espacial \textit{free-floating} de seis graus de liberdade, com foco em prevenção de colisão e desvio de obstáculos por meio de aprendizado por reforço. Os autores consideram explicitamente o efeito do acoplamento dinâmico entre a espaçonave e o manipulador, aspecto particularmente importante nesse tipo de sistema. Essa referência amplia a aplicação do aprendizado por reforço para além da aproximação translacional entre espaçonaves, alcançando também o problema de planejamento de movimento do manipulador em ambiente orbital.

De modo geral, a literatura revisada mostra que as pesquisas mais recentes caminham em direção a sistemas mais autônomos, integrados e orientados por restrições reais de missão. Enquanto trabalhos anteriores concentravam-se principalmente na formulação da dinâmica relativa e no controle de seguimento, as contribuições recentes passaram a enfatizar arquiteturas mais completas, que combinam percepção, planejamento, controle, prevenção de colisão e tomada de decisão assistida por técnicas de otimização e aprendizado. Esses aspectos são relevantes para a presente dissertação, pois mostram que a fase de aproximação e interação com o alvo deve ser compreendida não apenas como um problema dinâmico, mas também como parte de um processo operacional mais amplo, no qual segurança, autonomia e integração entre subsistemas assumem papel central.